% ===== PRÉAMBULE CANONIQUE — à coller VERBATIM en tête de CHAQUE .tex =====
\documentclass[11pt,a4paper]{article}
\usepackage[utf8]{inputenc}\usepackage[T1]{fontenc}\usepackage{lmodern}
\usepackage[french]{babel}
\usepackage{amsmath,amssymb,mathtools}\usepackage{icomma}
\usepackage[version=4]{mhchem}          % \ce{...} : équations & formules
\usepackage{siunitx}\sisetup{locale=FR} % \qty{}{}, \si{}, \ang{}
\usepackage{chemfig}                    % \chemfig{...} : structures développées
\usepackage{xcolor}\usepackage{colortbl}
\usepackage{tikz}\usepackage{pgfplots}\pgfplotsset{compat=1.18}
\usepackage[most]{tcolorbox}\usepackage{enumitem}\usepackage{array,booktabs}
\usepackage[a4paper,margin=2.2cm]{geometry}
\usetikzlibrary{arrows.meta,calc,angles,quotes,positioning,patterns,backgrounds,shapes.geometric}
\definecolor{primary}{RGB}{222,49,99}\definecolor{primarydark}{RGB}{150,22,55}
\definecolor{defcol}{RGB}{0,114,178}\definecolor{methcol}{RGB}{52,92,148}
\definecolor{excol}{RGB}{0,135,90}\definecolor{attncol}{RGB}{200,40,55}
\tcbset{boxbase/.style={enhanced,breakable,boxrule=0.9pt,arc=2.5pt,
  left=3mm,right=3mm,top=2.3mm,bottom=2.3mm,fonttitle=\bfseries,coltitle=white}}
% --- boîtes TUTORIEL (noms EXACTS, ne pas renommer) ---
\newtcolorbox{cle}[1][]{boxbase,colback=defcol!6,colframe=defcol,
  title={\textsc{À retenir}\ifx\relax#1\relax\relax\else\textmd{ : \itshape #1}\fi}}
\newtcolorbox{methode}[1][]{boxbase,colback=methcol!7,colframe=methcol,sharp corners,
  title={\textsc{Méthode}\ifx\relax#1\relax\relax\else\textmd{ --- \itshape #1}\fi}}
\newtcolorbox{exemplebox}[1][]{boxbase,colback=excol!5,colframe=excol,
  title={\textsc{Exemple}\ifx\relax#1\relax\relax\else\textmd{ : \itshape #1}\fi}}
\newtcolorbox{piege}[1][]{boxbase,colback=attncol!6,colframe=attncol,
  title={\textsc{Piège}\ifx\relax#1\relax\relax\else\textmd{ : \itshape #1}\fi}}
\newtcolorbox{remarque}[1][]{boxbase,colback=black!4,colframe=black!45,
  title={\textsc{Remarque}\ifx\relax#1\relax\relax\else\textmd{ : \itshape #1}\fi}}
% --- boîtes EXERCICE / CORRIGÉ (noms EXACTS, auto-comptées) ---
\newtcolorbox[auto counter]{exo}[1][]{boxbase,colback=primary!4,colframe=primary,
  title={\textsc{Exercice}~\thetcbcounter\ifx\relax#1\relax\relax\else\textmd{ --- \itshape #1}\fi}}
\newtcolorbox[auto counter]{corrige}[1][]{boxbase,colback=excol!4,colframe=excol,
  title={\textsc{Corrigé}~\thetcbcounter\ifx\relax#1\relax\relax\else\textmd{ --- \itshape #1}\fi}}
\newcommand{\R}{\mathbb{R}}\newcommand{\N}{\mathbb{N}}\newcommand{\Z}{\mathbb{Z}}\newcommand{\ds}{\displaystyle}
% (un \usepackage{fancyhdr}+en-tête est possible ; il est ignoré côté web)
% =========================================================================
\begin{document}
\begin{corrige}[reconnaître des isotopes]
On compare $Z$ (en bas) et $A$ (en haut).
\begin{enumerate}[label=\alph*)]
  \item \textbf{Oui} : même $Z=8$ (oxygène), $A$ différents ($16$ et $18$) : ce sont deux isotopes de
        l'oxygène.
  \item \textbf{Non} : $Z$ différents ($18$ contre $20$), donc ce sont deux \emph{éléments} différents
        (argon et calcium), même s'ils ont le même $A=40$.
  \item \textbf{Oui} : même $Z=6$ (carbone), $A$ différents ($12$ et $14$) : deux isotopes du carbone.
\end{enumerate}
\end{corrige}

\begin{corrige}[neutrons et chimie des isotopes]
\begin{enumerate}[label=\alph*)]
  \item $n(n^0) = A - Z$ : $\ce{^{24}Mg}$ a $12$, $\ce{^{25}Mg}$ a $13$, $\ce{^{26}Mg}$ a $14$ neutrons.
  \item À l'état neutre, $n(e^-) = Z = 12$ électrons pour les trois.
  \item Oui : ayant le même $Z$, ils ont le même cortège électronique ($12$ électrons, même
        configuration), donc \emph{les mêmes propriétés chimiques} ; ils ne diffèrent que par leur masse.
\end{enumerate}
\end{corrige}

\begin{corrige}[masse atomique moyenne à deux isotopes]
On applique $\overline{A} = p_1 A_1 + p_2 A_2$ (abondances en fractions).
\begin{enumerate}[label=\alph*)]
  \item $\overline{A} = 0{,}75\times 35 + 0{,}25\times 37 = 26{,}25 + 9{,}25 = 35{,}5$.
  \item $\overline{A} = 0{,}70\times 63 + 0{,}30\times 65 = 44{,}1 + 19{,}5 = 63{,}6$.
  \item $\overline{A} = 0{,}52\times 107 + 0{,}48\times 109 = 55{,}64 + 52{,}32 = 107{,}96$.
\end{enumerate}
\end{corrige}

\begin{corrige}[retrouver une abondance manquante]
On utilise $\sum p_i = 100\,\%$.
\begin{enumerate}[label=\alph*)]
  \item $100 - 60 = 40\,\%$.
  \item $100 - (78 + 10) = 100 - 88 = 12\,\%$.
  \item $100 - (5{,}8 + 91{,}7 + 2{,}2) = 100 - 99{,}7 = 0{,}3\,\%$.
\end{enumerate}
\end{corrige}

\begin{corrige}[l'unité de masse atomique]
\begin{enumerate}[label=\alph*)]
  \item $\SI{1}{uma} \approx \SI{1,66054e-27}{\kilo\gram}$ (par définition, $\tfrac{1}{12}$ de la masse
        d'un atome de $\ce{^{12}C}$).
  \item $\ce{^{23}Na}$ compte $23$ nucléons, chacun $\approx 1$ uma : sa masse est $\approx 23$ uma.
  \item $\SI{1}{uma} = \dfrac{1}{N_{\!A}}\,\si{\gram} = \dfrac{1}{\SI{6,022e23}{}}\ \si{\gram}
        \approx \SI{1,66054e-24}{\gram}$.
\end{enumerate}
\end{corrige}

\end{document}
